\documentclass[11pt,a4paper]{article}

% ==================================================
% LINGUA E CODIFICA
% ==================================================
\usepackage[T1]{fontenc}
\usepackage[utf8]{inputenc}
\usepackage[italian]{babel}

% ==================================================
% FONT
% ==================================================
\usepackage{lmodern}
\usepackage{microtype}

% ==================================================
% GEOMETRIA DELLA PAGINA
% ==================================================
\usepackage[
a4paper,
top=2.2cm,
bottom=2.2cm,
left=2.5cm,
right=2.5cm
]{geometry}

% ==================================================
% MATEMATICA
% ==================================================
\usepackage{amsmath}
\usepackage{amssymb}
\usepackage{amsthm}
\usepackage{mathtools}
\usepackage{bm}

% ==================================================
% GRAFICA
% ==================================================
\usepackage{graphicx}
\usepackage{float}
\usepackage{caption}
\usepackage{subcaption}

% ==================================================
% TABELLE
% ==================================================
\usepackage{booktabs}
\usepackage{array}
\usepackage{tabularx}
\usepackage{multirow}

% ==================================================
% COLORI
% ==================================================
\usepackage{xcolor}

\definecolor{myblue}{RGB}{30,90,160}
\definecolor{mygreen}{RGB}{40,130,90}
\definecolor{myred}{RGB}{180,60,60}
\definecolor{lightgray}{RGB}{245,245,245}

% ==================================================
% LINK
% ==================================================
\usepackage[
colorlinks=true,
linkcolor=myblue,
urlcolor=myblue,
citecolor=mygreen
]{hyperref}

% ==================================================
% HEADER E FOOTER
% ==================================================
\usepackage{fancyhdr}

\pagestyle{fancy}
\fancyhf{}

\lhead{\textit{Geometria}}
\rhead{\textit{Appunti}}
\cfoot{\thepage}

\renewcommand{\headrulewidth}{0.4pt}

% ==================================================
% TITOLI DELLE SEZIONI
% ==================================================
\usepackage{titlesec}

\titleformat{\section}
{\Large\bfseries\color{myblue}}
{\thesection}
{1em}
{}

\titleformat{\subsection}
{\large\bfseries}
{\thesubsection}
{1em}
{}

\titleformat{\subsubsection}
{\normalsize\bfseries}
{\thesubsubsection}
{1em}
{}

% ==================================================
% SPAZIATURA
% ==================================================
\usepackage{setspace}

\setstretch{1.1}

\setlength{\parindent}{0pt}
\setlength{\parskip}{6pt}

% ==================================================
% LISTE
% ==================================================
\usepackage{enumitem}

\setlist[itemize]{
	itemsep=2pt,
	topsep=4pt
}

\setlist[enumerate]{
	itemsep=2pt,
	topsep=4pt
}

% ==================================================
% BOX
% ==================================================
\usepackage[most]{tcolorbox}

\newtcolorbox{definitionbox}{
	colback=blue!4,
	colframe=myblue,
	boxrule=0.8pt,
	arc=2mm,
	title=Definizione,
	fonttitle=\bfseries
}

\newtcolorbox{importantbox}{
	colback=red!4,
	colframe=myred,
	boxrule=0.8pt,
	arc=2mm,
	title=Importante,
	fonttitle=\bfseries
}

\newtcolorbox{examplebox}{
	colback=green!4,
	colframe=mygreen,
	boxrule=0.8pt,
	arc=2mm,
	title=Esempio,
	fonttitle=\bfseries
}

% ==================================================
% TEOREMI
% ==================================================

\theoremstyle{plain}

\theoremstyle{remark}

\newtheorem*{remark}{Osservazione}

\theoremstyle{definition}
\newtheorem{propmanuale}{Proposizione}
\newenvironment{proposizione}[1]
{\renewcommand{\thepropmanuale}{#1}\propmanuale}
{\endpropmanuale}

\theoremstyle{definition}
\newtheorem{defmanuale}{Definizione}
\newenvironment{definizione}[1]
{\renewcommand{\thedefmanuale}{#1}\defmanuale}
{\enddefmanuale}

\theoremstyle{plain}
\newtheorem{teoremainterno}{Teorema}
\newenvironment{teorema}[1]
{\renewcommand{\theteoremainterno}{#1}\begin{teoremainterno}}
	{\end{teoremainterno}}

% ==================================================
% CODICE
% ==================================================
\usepackage{listings}

\lstset{
	backgroundcolor=\color{lightgray},
	basicstyle=\ttfamily\small,
	frame=single,
	breaklines=true,
	numbers=left,
	numberstyle=\tiny,
	tabsize=4,
	showstringspaces=false
}

% ==================================================
% COMANDI MATEMATICI UTILI
% ==================================================

\newcommand{\R}{\mathbb{R}}
\newcommand{\N}{\mathbb{N}}
\newcommand{\Z}{\mathbb{Z}}
\newcommand{\Q}{\mathbb{Q}}
\newcommand{\C}{\mathbb{C}}

\newcommand{\dd}{\,\mathrm{d}}
\newcommand{\e}{\mathrm{e}}
\newcommand{\ii}{\mathrm{i}}

\DeclareMathOperator{\rank}{rank}
\DeclareMathOperator{\Ker}{Ker}
\DeclareMathOperator{\Img}{Im}
\DeclareMathOperator{\Span}{span}

% ==================================================
% DOCUMENTO
% ==================================================

\begin{document}
	
	\begin{center}
		
		{\LARGE\bfseries Geometria}
		
		\vspace{0.3cm}
		
		{\large Appunti}
		
		\vspace{0.2cm}
		
	\end{center}
	
	\vspace{0.5cm}
	
	\tableofcontents
	
	\newpage
	
	% ==================================================
	
	\section{Introduzione}
	
	Questi sono appunti presi dal libro "Casali, Gagliardi, Grasselli" e scritti in linguaggio più simile ad Algebra Astratta (leggasi: più formale). Essi sono stati redatti al fine di studiare il corso di Geometria I del CdL "Informatica" (L-31) dell'Università di Napoli.
	
	\section{Strutture Algebriche}
	Per struttura algebrica intendiamo una n-pla costituita da insiemi e operazioni su di essi.
	Un esempio è $(X, \perp)$, dove $X$ è un insieme e $\perp$ è un'operazione binaria interna a $X$.
	
	\subsection{Gruppi}
	Sia $(X, \perp)$ un gruppoide.
	
	\begin{enumerate}[label=\roman*)]
		\item L'operazione $\perp$ è detta associativa se
		$$\forall x,y,z \in X, (x \perp y) \perp z = z \perp (y \perp z)$$
		
		\item L'operazione $\perp$ è detta commutativa se
		$$ \forall x,y \in X, x \perp y = y \perp x $$
		
		\item Un elemento $u \in X$ è detto elemento neutro per l'operazione $\perp$ se
		$$ \forall x \in X, x \perp u = u \perp x = x $$ 
		
		\item Se $(X, \perp)$ ammette elemento neutro $u$, un elemento $x \in X$ è detto invertibile se
		$$ \exists x' \in X, x \perp x' = x' \perp x = u $$
	\end{enumerate}
	
	\begin{proposizione}{2.3}
		 \leavevmode\vspace{-0.5\baselineskip}
		\begin{enumerate}[label=\roman*)]
			
			\item Se $(X, \perp)$ ammette elemento neutro, questo è unico
			
			\item Se $(X, \perp)$ ammette elemento neutro e $\perp$ è associativa, allora $x \in X$ è invertibile e il suo inverso è unico
			
			\item Se $(X, \perp)$ ammette elemento neutro e $\perp$ è associativa e se $x,y \in X$ sono invertibili, allora anche $x \perp y$ è invertibile. Inoltre, detto $(x \perp y)'$ l'inverso di $(x \perp y)$ si ha che $(x \perp y)' = y' \perp x'$
		\end{enumerate}
	\end{proposizione}
	
	\begin{proof}
		\leavevmode\vspace{-0.5\baselineskip}
		\begin{enumerate}[label=\roman*)]
			
			\item Siano $u$ e $\bar{u}$ due elementi neutri per $(X, \perp)$. Avremo che:
			$$ u = u \perp \bar{u} = \bar{u} $$
			Dunque necessariamente devono coincidere
			
			\item Sia $u$ l'elemento neutro di di $(X, \perp)$ e siano $x', x''$ entrambi inversi di $x$
			$$x' = x \perp u = x' \perp (x \perp x'') = (x' \perp x) \perp x'' = u \perp x'' = x''$$
			
			\item Dimostriamo dunque che l'inverso di $(x \perp y)$ esiste ed è $y' \perp x'$. Per dimostrare che $x'$ è inverso di $x$ dobbiamo dimostrare che $x \perp x' = x' \perp x = u$.
			\\
			\textit{Lato destro:}
			$$ (x \perp y) \perp (y' \perp x') = \underbrace{x \perp (y \perp y') \perp x'}_{\text{Associatività}} = x \perp u \perp x' = x \perp x' = u $$
			\textit{Lato sinistro:}
			$$(y' \perp x') \perp (x \perp y) = y' \perp (x' \perp x) \perp y = y' \perp u \perp y = y' \perp y = u$$
		\end{enumerate}
	\end{proof}
	
	\begin{definizione}{2.1}
		Sia $G$ un insieme e $\perp$ un'operazione binaria interna su G. Diremo che la struttura algebrica $(G, \perp)$ è un gruppo se soddisfa le seguenti proprietà:
		
		\begin{enumerate}[label=\roman*)]
			
			\item $\perp$ è associativa
			\item Ammette elemento neutro
			\item Ogni elemento di G è invertibile
			
		\end{enumerate}
		
		Un gruppo è detto abeliano (o commutativo) se soddisfa un'ulteriore proprietà:
		
		\begin{enumerate}[label=\roman*), start=4]
			\item $\perp$ è commutativa
		\end{enumerate}
		
	\end{definizione}
	
	\subsection{Anelli e campi}
	
	\begin{definizione}{2.5 (Anello)}
		Sia $A$ un insieme e $+$ e $\times$ due operazioni interne su $A$. Diremo che la struttura algebrica $(A, +, \times)$ è un anello se soddisfa le seguenti proprietà:
		
		\begin{enumerate}[label=\roman*)]
			\item $(A, +)$ è un gruppo abeliano
			\item $\times$ è associativo
			\item $\times$ è distributivo rispetto a $+$, ovverosia:
			$$x \times (y+z) = (x \times y) + (x \times z)$$
			\begin{center}
				e
			\end{center}
			$$(y+z)\times x = (y \times x) + (z \times x)$$
		\end{enumerate}
		
		Un anello $(A, +, \times)$ è detto commutativo se $\times$ è commutativo ed è detto unitario se $(A, \times)$ ammette elemento neutro.
		
	\end{definizione}
	
	\begin{definizione}{2.6 (Campo)}
		$(K, +, \times)$ è detto campo se è un anello commutativo unitario in cui ogni elemento di $(K, +)$ eccetto l'elemento neutro è invertibile rispetto a $\times$
		
		In altre parole:
		\begin{enumerate}[label=\roman*)]
			\item $(K, +)$ è un gruppo abeliano
			\item Sia $u_{+}$ l'elemento neutro di $(K, +)$. Allora $(K - \{ u_{+} \}, \times)$ è un gruppo abeliano
			\item $\times$ è distributivo rispetto a $+$
		\end{enumerate}
		
	\end{definizione}
	
	\textbf{Osservazione:} Notiamo che ad un primo sguardo, le due definizioni non coinciderebbero. Infatti affinché $(K, +, \times)$ sia un anello commutativo unitario servono cinque cose: che $(K, +)$ sia un gruppo abeliano, che $\times$ sia distributivo rispetto a $+$, e che $\times$ sia associativo, commutativo e che abbia elemento neutro su tutto $K$. Le prime due sono i punti \textit{i)} e \textit{iii)}, mentre per quanto riguarda le altre tre, la definizione \textit{ii)} ci dice invece che queste proprietà valgono su $K - \{u_{+}\}$, non su $K$. Come conciliare le due cose? Vediamo: \\
	\textit{Commutatività su tutto $K$:} Siano $a,b \in K$. Se $a$ e $b$ sono entrambe diverse da $0$, allora vale che $a \times b = b \times a$ grazie al punto \textit{ii)}. Se una delle due è $0$, allora vale che $a \times 0 = 0 \times a = 0$ come verrà dimostrato nel Teorema 1.1 \\
	\textit{Associatività su tutto $K$:} Siano $a,b,c \in K$. Se sono tutti diversi da $0$ allora vale che $(a \times b) \times c = a \times (b \times c)$ grazie al punto \textit{ii)}. Se almeno uno è $0$ allora $(a \times b) \times c = 0 = a \times (b \times c)$ grazie al risultato che verrà mostrato nel Teorema 1.1\\
	\textit{Elemento neutro su tutto K:} Sia $u_{\times}$ l'elemento neutro di $(K - \{u_{+}\}, \times)$. Banalmente, è anche elemento neutro di $K$ in quanto basta mostrare che $u_{+} \times u_{\times} = u_{+}$
	
	\textbf{Osservazione:} Nella definizione di campo, escludiamo l'elemento neutro di $(K, +)$ dall'imposizione di essere invertibile rispetto a $\times$ perché $u_{+}$ non è invertibile in un anello non banale. Infatti sia $u_{+}'$ l'inverso di $u_{+}$. Dovrebbe valere che $u_{+}' \times u_{+} = u_{\times}$.
	In realtà accade che $\forall a \in K, a \times u_{+} = u_{+}$
	
	\begin{teorema}{1.1}
		Sia $(A, +, \times)$ una struttura algebrica e sia $a \in A$. Sia $(A, +)$ un gruppo con elemento neutro $u_{+}$ e che $\times$ sia distributiva su $+$. Allora $a \times u_{+} = u_{+}$.
		
		\begin{proof}
			\begin{align*}
				&a \times u_{+} \\ &= a \times (u_{+} + u_{+}) \\
				&= (a \times u_{+}) + (a \times u_{+}) \\
				& \Rightarrow
				 (a \times u_{+}) - (a \times u_{+}) = ((a \times u_{+}) + (a \times u_{+})) - (a \times u_{+}) \\
				& \Rightarrow
				 u_{+} = (a \times u_{+}) + ((a \times u_{+}) - (a \times u_{+})) \\
				& \Rightarrow
				 u_{+} = (a \times u_{+}) + u_{+} \\
				& \Rightarrow
				 u_{+} = (a \times u_{+})
			\end{align*}
		\end{proof}
		
		Vale anche il contrario, ovvero che $u_{+} \times a = u_{+}$ che si dimostra nello stesso modo
		
	\end{teorema}
	
	\textbf{Osservazione:} Ogni campo ha sempre perlomeno due elementi, ovvero i due elementi neutri delle due operazioni. Questo perché
	\begin{enumerate}[label=\roman*)]
		\item $(K, +)$ è un gruppo abeliano, dunque ammette elemento neutro $u_{+}$
		\item $(K - \{u_{+}\}, \times)$ è un gruppo abeliano, dunque ammette elemento neutro, che sarà necessariamente diverso da $u_{+}$ in quanto escluso dall'insieme
	\end{enumerate}
	
	D'ora in poi $u_{+}$ verrà indicato anche con $0$ e $u_{\times}$ verrà indicato anche con $1$ per allinearli agli elementi neutri dell'esempio classico di campo: $(\mathbb{R}, +, \cdot)$
	
	\begin{definizione}{2.9}
		Sia $(A, +, \cdot)$ un anello. Se $a,b \in A - \{0\}$ e $a \cdot b = 0$ si dice che $a$ e $b$ sono divisori dello zero
	\end{definizione}
	
	
	\textbf{Osservazione 2.10}
	La mancanza di divisori dello zero nell'anello $A$ equivale alla seguente proprietà, nota come "legge di annullamento del prodotto":
	$$\forall a,b \in A, a \cdot b = 0\Rightarrow a = 0 \text{ oppure } b=0$$
	
	\begin{proposizione}{2.11}
		Ogni campo è privo di divisori dello zero
		\begin{proof}
			Sia $(K, +, \cdot)$ un campo e siano $a,b \in K$ tali che $a \cdot b = 0$. Se $a \neq 0$ allora:
			$$b = (a^{-1} \cdot a) \cdot b = a^{-1} \cdot (a \cdot b) = a^{-1} \cdot 0 = 0$$
			In altre parole abbiamo dimostrato che se uno degli elementi è diverso da $0$, allora l'altro deve necessariamente essere uguale a $0$. \textit{(Nota: $a^{-1}$ è l'inverso di $a$ per $(K, \cdot)$)}
		\end{proof}
	\end{proposizione}
	
	\newpage
	\section{Spazi e sottospazi vettoriali}
	\subsection{Spazi vettoriali}
	Sia V un insieme e $(K, +, \cdot)$ un campo. 
	\begin{definizione}{4.1}
		Diremo spazio vettoriale su K una struttura algebrica $(K, V, \boxplus, \boxdot)$ costituita da un campo $K$ i cui elementi sono detti scalari e da un insieme $V$ i cui elementi sono detti vettori, da un'operazione binaria interna su V detta "somma di vettori":
		$$\boxplus: V \times V \rightarrow V$$
		e da un'operazione detta "prodotto di uno scalare per un vettore":
		$$\boxdot: K \times V \rightarrow V$$
		
		Inoltre deve soddisfare questi assiomi $\forall \alpha, \beta \in K, \forall v, w \in V$:
		\begin{enumerate}
			\item $(V, \boxplus)$ è un gruppo abeliano
			\item $(\alpha + \beta) \boxdot v = (\alpha \boxdot v) \boxplus (\beta \boxdot v)$
			\item $\alpha \boxdot (v \boxplus w) = (\alpha \boxdot v) \boxplus (\alpha \boxdot w)$
			\item $\alpha \boxdot (\beta \boxdot v) = (\alpha \cdot \beta) \boxdot v$
			\item $1 \boxdot v = v$
		\end{enumerate}
	\end{definizione}
	
	L'insieme $V$ è detto sostegno dello spazio vettoriale. Inoltre, per leggerezza di notazione, $\boxdot \text{ e } \boxplus$ verranno d'ora in poi indicati con gli stessi simboli che denotano le operazioni in $K$, ovvero, rispettivamente, $\cdot \text{ e } +$. Tale abuso di notazione non sarà ambiguo perché sarà possibile capire il vero tipo di operazione dal tipo di operandi a cui è applicata.
	Inoltre, sempre per leggerezza di notazione, indicheremo uno spazio vettoriale con il solo nome del suo sostegno.
	
	\begin{proposizione}{4.2}
		Sia $V$ uno spazio vettoriale sul campo $K$. Sia $0_{v}$ l'elemento neutro in $(V, \boxplus)$. Per ogni $\alpha, \beta \in K, \forall v,u \in V$ si ha:
		\begin{enumerate}
			\item $\alpha v = 0_{v} \iff \alpha = 0 \text{ oppure } v = 0_{v}$
			\item $(-\alpha)v = \alpha(-v) = -(\alpha v)$
			\item Se $\alpha v = \beta v \text{ e } v \neq 0 \Rightarrow \alpha = \beta$
			\item $\alpha u = \alpha v \text{ e } \alpha \neq 0 \Rightarrow u = v$
		\end{enumerate}
		\begin{proof}
			\begin{enumerate}
				\item
				$\Leftarrow$
			\end{enumerate}
		\end{proof}
	\end{proposizione}
	
\end{document}